%============================================================ 2. 관련 연구
\section{관련 연구}
\label{sec:related}

\subsection{무기-표적 할당과 조합 최적화}

방어 자산을 표적에 배분하는 문제는 무기-표적 할당(WTA)으로 정식화되어 정확해법과
휴리스틱이 모두 연구되었다~\citep{ahuja2007wta}. 1:1 배정으로 축소하면 헝가리안
알고리즘~\citep{kuhn1955hungarian}으로 다항 시간에 최소비용 매칭을 얻을 수 있다.
이러한 방법은 비용 함수가 명시적으로 주어질 때 최적을 보장하지만, 비용 함수 자체가 전술적
의도를 담아야 한다는 부담을 남긴다. ``이 무리는 양동이므로 무시하고 뒤쪽을 대비하라''와
같은 판단은 거리·각도 비용으로 표현되지 않는다.

\emph{본 연구와의 차이}: 본 연구는 조합 최적화를 배정 계층에서 제거하고 그 자리를 언어모델의
판단으로 대체한다. 이는 최적성을 포기하는 대신 비용 함수로 표현되지 않는 전술적 맥락을
반영하기 위한 선택이며, 그 대가를 \S\ref{sec:results}에서 정량적으로 보고한다.

\subsection{무인 위협 대응}

대무인기 체계에 관한 조사~\citep{wang2021cuas}는 대응을 탐지·식별·무력화로 나누고,
무력화를 다시 연성(전자전)·경성(물리 타격)·포획(그물 등)으로 구분한다. 유선 조종 표적의
확산~\citep{sutea2026fiber}은 연성 대응의 적용 범위를 좁혔고, 요격 비용의 비대칭
~\citep{rumbaugh2024cost}은 경성 대응의 지속성을 제약한다. 해상 영역에서 다중 USV의 협력
기동은 오래 연구되어 왔으며~\citep{liu2016usv}, 최근에는 다중 에이전트 강화학습으로 회피
표적을 협조 포위하는 정책이 실선 검증과 함께 보고되었다~\citep{zhang2025encircle}.

\emph{본 연구와의 차이}: 선행 연구의 포위(encirclement)는 방어정이 표적보다 빠르거나 대등한
기동 성능을 전제한다. 본 연구는 방어정이 \emph{더 느린} 조건에서 추격 대신 접근 회랑의
선제 차단을 다루며, 표적을 따라가는 것이 아니라 표적이 지나갈 자리에 그물벽을 놓는 문제로
정식화한다.

\subsection{다중 에이전트 강화학습과 크레딧 배분}

중앙 집중 학습·분산 실행(CTDE) 계열에서 MADDPG~\citep{lowe2017maddpg}는 중앙 비평가를,
QMIX~\citep{rashid2018qmix}는 단조 혼합망으로 팀 가치를 개체 가치로 분해하는 구조를
제시하였다. COMA~\citep{foerster2018coma}는 한 에이전트의 행동만 주변화한 반사실적
기준선으로 개별 기여를 분리한다. MAPPO~\citep{yu2022mappo}는 PPO~\citep{schulman2017ppo}를
다중 에이전트로 확장하여 강한 기준선을 제공한다. 이들은 모두 가치망(비평가)을 학습한다.
본 논문이 이 계열을 실험 기준선으로 쓰지 \emph{않은} 이유는 \S\ref{sec:limitations}에 밝힌다
--- 행동 표현이 달라져 표현과 알고리즘의 기여가 뒤섞이기 때문이다.

가치망을 두는 구조는 강화학습의 표준적 선택이지만~\citep{sutton2018rl}, 가치 추정의 편향과
학습 불안정이 성능을 좌우하는 요인이 된다. 연속 제어에서는 엔트로피 정칙화로 이를 완화하는
접근도 널리 쓰인다~\citep{haarnoja2018sac}. 한편 언어모델 정렬 분야에서 제안된 그룹 상대
정책 최적화(GRPO)~\citep{shao2024deepseekmath}는 같은 상태에서 뽑은 여러 후보의 보상을 서로
비교해 이득을 산정함으로써 가치망 자체를 제거한다.

\emph{본 연구와의 차이}: 본 연구는 GRPO를 다중 에이전트 제어로 옮기고, COMA의 반사실적
착상을 \emph{학습된 비평가 없이} 시뮬레이터 재진행으로 직접 구현한다. 즉 한 척의 행동만
바꾸어 실제로 굴려 본 차이를 그 배의 몫으로 돌린다. 가치 추정 편향이 원천적으로 없고,
비평가 학습의 불안정성도 발생하지 않는다.

\subsection{지도 표현과 포인터 방식 행동}

의미론적 분할은 입력 해상도의 조밀 예측을 가능하게 하였다 --- 완전 합성곱
구조~\citep{long2015fcn}와 U-Net~\citep{ronneberger2015unet}이 그 계보다. 합성곱은 본질적으로 평행이동 등변이라 절대
위치를 흐리는데, CoordConv~\citep{liu2018coordconv}는 좌표 채널을 덧붙여 이를 보완한다.
FiLM~\citep{perez2018film}은 조건 벡터로부터 채널별 배율과 이동값을 만들어 특징 지도를
변조하는 경량 조건화 기법이다. 한편 이산 후보 위에서 원소를 \emph{지목}하는 포인터
방식~\citep{vinyals2015pointer}은 출력 어휘가 입력에 의존하는 문제에 적합하며, 어텐션
연산~\citep{vaswani2017attention}의 점수 계산과 같은 형식을 공유한다.

\emph{본 연구와의 차이}: 선행 연구에서 분할 구조는 지각(무엇이 어디에 있는가)에, 포인터는
조합 최적화(어느 원소를 고를 것인가)에 쓰였다. 본 연구는 둘을 결합하여, 분할 구조가 낸
조밀 점수맵 위에서 \emph{행동}으로서의 픽셀을 지목한다. 즉 점수맵이 지각 결과가 아니라
정책의 행동 분포 자체가 된다. 이로써 행동공간이 격자로 이산화되어 규제되고, 마스킹으로
안전 제약(육지·사거리·배정 회랑)을 행동 분포에 직접 부과할 수 있다.

\subsection{언어모델 기반 계획}

언어모델을 계획 주체로 사용하는 연구는 고수준 지시를 실행 가능한 단계로 분해하는 방향으로
발전해 왔다~\citep{huang2022zeroshot}. SayCan~\citep{ahn2022saycan}은 언어모델의 사전 지식과
로봇의 실행 가능성 추정을 결합하여 실행 불가능한 계획을 걸러낸다. 출력 형식을 스키마로
강제하는 구조적 출력~\citep{openai2026function}은 언어모델 출력을 하위 시스템이 파싱 가능한
형태로 고정한다.

\emph{본 연구와의 차이}: 선행 연구는 대체로 언어모델의 계획을 실행 가능성 필터나 최적화로
보정한다. 본 연구는 그 보정을 제거한 조건을 \emph{주 실험 조건}으로 삼아, 보정 없이 언어모델
판단만으로 성립하는지, 그리고 어느 능력 수준부터 이득이 되는지를 측정한다.
